\chapter{Structural adaptation and lifecycle governance}
\label{evod43:chapter}

A candidate improves a quality score from .80 to .95 but lowers a safety
invariant score from 1 to .5. Another reaches quality .85 while retaining
safety 1. Which may become active? The first fails the required gate despite
its larger quality gain. A system that can propose new structures is not
thereby entitled to deploy them.

Chapter 13 distinguished transient state, learned weights, structural proposals
and population selection. Chapters 39--41 supplied typed operations, memory
and legal compilation. This chapter governs changes to those objects through
immutable identity, evidence, bounded canaries, promotion and rollback.
Its local registry is a serial specification, not a signed release service or
a demonstration of autonomous scientific improvement.

\section{Name the object that changes}
\label{evod43:objects}
\index{adaptation!structural versus parametric}\index{artifact!immutable lineage}
Candidates can change weights, adapters, memory policies, retrieval indexes,
compiler rules, prompts/policies or topology. These have different evaluation
and state-compatibility requirements. A token-level router selecting an expert
from a fixed graph is not automatically structural learning. Updating recurrent
activation state is not parameter learning. ``Self-improvement'' is too coarse
a type for an engineering change.

Let $A$ be an immutable artifact and $T$ a declared transformation with recipe,
data and budget. A proposal creates $A'=T(A)$ with explicit parent identity;
it never rewrites the active artifact's bytes. The executable registry stores
payload, recipe, training-data identity and a native-state fingerprint in the
content identity. Lifecycle stage and measured metrics live in separate records.
Changing a stage therefore does not change artifact bytes or its identity.

\Listing{Registry.create}
Duplicate creation is idempotent: it returns the existing artifact without
resetting a rejected or active stage. A parent must already exist. With a
fresh child identity and immutable parent field, every new parent edge points
to an older registered object, giving an acyclic lineage. Multiple children
can share a parent, preserving rejected branches as evidence rather than
rewriting the history around a winning candidate.

\Plate{01-lineage}{The unsafe branch remains in lineage but never serves. Promotion and rollback move serving identity; they do not reverse parent edges.}{evod43:lineage}

The hash binds these declared strings; it does not prove that named data were
actually used, that the payload is a real trained model or that a producer is
authorized. Production packaging must bind verified artifact bytes and complete
configuration, and authenticated release mechanisms must enforce authority.
Here role labels are trusted caller inputs and strings are explicit fixture
provenance, not cryptographic attestations.

\section{A gate is a conjunction of constraints}
\label{evod43:gates}
\index{evaluation!conjunctive gates}\index{nonfinite metrics}
For a metric vector $m$, define each gate as a metric name, direction and
threshold. Some metrics require lower bounds, others upper bounds:
\[
 Eligible(m)=\bigwedge_{i\in I_{min}}[m_i\ge g_i]
 \ \wedge\!\bigwedge_{j\in I_{max}}[m_j\le g_j].
\]
Our fixed registry policy requires quality at least .8, safety at least 1,
and latency at most 10. Latency uses fixture milliseconds; quality and safety
are stipulated dimensionless scores. They are not measured validation accuracy
or a general safety probability. Missing metrics fail a gate; NaN, infinity,
Boolean values and invalid gate directions are rejected before stage mutation.
An empty gate set is not accepted through vacuous truth.

\Listing{gates_pass}
Gates belong to the registry's configured policy, not to the proposing candidate.
The evaluator cannot pass a new empty or weaker gate set on each call.
More realistic policies include confidence intervals, absolute invariants,
regression budgets and workload slices. Hard constraints such as tenant
isolation, authority-before-side-effect, bounded memory and complete lifecycle
traces cannot be traded for a higher aggregate task score.

If all $k$ deterministic gates pass, eligibility is true by conjunction.
That logical statement says nothing about estimation error. With statistical
metrics, each reported estimate needs its sample size and uncertainty. Testing
many candidates against the same held-out set also adapts to that set;
selection then needs a fresh final evaluation or an explicit sequential/multiple
testing design. Calling the selection set held-out does not make it untouched.

\section{Bind evaluation to candidate, protocol and budget}
\label{evod43:evaluation}
\index{evaluation!held-out data}\index{population search!budgets}
Evaluation is legal once from candidate stage. Its immutable report binds
artifact ID, suite version, evaluation-data identity, evaluator, budget,
metric tuple and pass decision. The registry requires its configured held-out
data/suite labels and rejects data equal to the artifact's declared training
data. A second evaluation of the same object is rejected; a new versioned
candidate/evidence protocol is needed rather than overwriting its history.

These checks prevent certain bookkeeping errors, not data leakage by a dishonest
producer. Dataset labels are not content hashes or proof of non-overlap.
A real evaluation artifact must record exact examples/splits, task conventions,
code, seeds, compute budgets, hardware, aggregation and uncertainty. Genomic
splits may need sequence/homology/species controls; tool safety needs negative
and adversarial workloads. This registry executes none of those experiments.

Population search should compare siblings under matched budgets or report
their differences. Preserve failed candidates, negative controls and the
selection rule. The learned or search-generated proposal can be wrong while
the governance mechanism correctly rejects it. Conversely, a passed synthetic
registry test does not establish that a structural change is useful.

\section{Separate parametric stability from structural legality}
\label{evod43:ewc}
\index{elastic weight consolidation}\index{stability and plasticity}
EWC provides a concrete parametric stability/plasticity method. Its paper
uses a diagonal Fisher-weighted quadratic around previous parameters,
\[
 L(\theta)=L_B(\theta)+\frac\lambda2
       \sum_i F_i(\theta_i-\theta^*_{A,i})^2,
\]
motivated by a local approximate posterior~\cite{evod43:ewc}.
This is an approximate method, not a universal guarantee against forgetting
and not an authorization protocol. The empirical tasks in that paper are not
reproduced by our registry.

To see the tradeoff without claiming trained behavior, stipulate one-coordinate
task loss $L_B=\tfrac c2(\theta-b)^2$, previous optimum $a$, importance $F\ge0$,
$c>0$ and $\lambda\ge0$. Differentiate:
\[
 \frac{dL}{d\theta}=c(\theta-b)+\lambda F(\theta-a)=0,
 \qquad
 \theta^*=\frac{cb+\lambda Fa}{c+\lambda F}.
\]
The second derivative $c+\lambda F>0$ proves a unique global minimizer for
this stipulated quadratic. With $a=0,b=2,c=1,F=3$, it becomes
$2/(1+3\lambda)$. Increasing the penalty keeps the coordinate nearer the
old optimum while sacrificing the new task optimum.
\begin{center}\input{\Generated/quadratic.tex}\end{center}
Tests compare this closed form with an independent finite-difference derivative
and endpoint objective values. The numbers are algebra, not a neural-network
experiment or an estimated Fisher matrix.

Structural changes require different checks. Adding a memory operator,
rerouting a tool or changing a compiler rule can alter dependency topology
and effect order even if weights are unchanged. Chapter 41's original/optimized
observations, verifier failures and native interfaces need regression coverage.
A parameter regularizer cannot establish those compiler properties.

\section{Make lifecycle transitions explicit}
\label{evod43:lifecycle}
\index{canary!statistical limits}\index{shadow evaluation}
The stages are candidate, evaluated, canary, active, retired and rejected.
Passing offline gates moves candidate to evaluated; failure moves it to
rejected. A bounded successful canary is required before promotion.
Rejected objects have no direct serving path. Changing a failed object and
trying again creates a new immutable child with new evidence.

\Plate{02-lifecycle}{Guards constrain every serving path. Rollback has its own eligibility rule and cannot bypass rejection by naming an arbitrary object.}{evod43:stages}

Shadow evaluation observes production-like inputs without controlling visible
outputs. It must still isolate memory writes and tools; a shadow call that
executes real effects is not harmless. Canary deployment serves a bounded
fraction under predefined limits and rollback conditions. It follows offline
correctness gates rather than using user exposure to discover basic type or
authorization failures.

Our canary specification accepts an exposure fraction in $(0,.1]$, at least
20 observations, valid error counts and nonnegative p95 latency. Zero errors
and p95 at most 10 move the candidate to canary; otherwise it is rejected.
These are caller-supplied observations, not actual traffic routing or quantile
measurement. A drain flag later is also a caller assertion, not implemented
request draining.

Even honest zero errors among 20 independent Bernoulli trials gives weak
statistical assurance. If the true error probability is $p$, the chance of
zero errors is $(1-p)^{20}$. Solving $(1-p)^{20}=.05$ gives one-sided upper
bound $p=1-.05^{1/20}\approx .139$. A small zero-error canary is therefore a
transition fixture, not evidence of a tiny production failure rate. Dependence,
nonstationarity and workload slices further limit that calculation.

\section{Promote with authority and a current revision}
\label{evod43:promotion}
\index{Promotion!revision guard}\index{authority!release roles}
Proposal, evaluation and promotion have separate caller roles in the reference:
search, lab and release. Promotion checks passed offline/canary evidence and
the current serving revision before changing any stage or pointer. Two
release decisions based on revision $r$ cannot both succeed against the
serialized store: the first increments it, making the second stale.

\Listing{Registry._switch}
The switch preflights all checks before mutating active identity, stages,
history, revision and audit tuple. It retires the old active object, activates
the target and appends a transition. This is one serialized method under
trusted callers. It is not a database transaction: an interpreter crash,
allocation failure or concurrent unsynchronized caller can break assumptions.
A real implementation needs a durable transaction/CAS over the whole state
and audit record, authenticated role separation and verified signatures/hashes.
Chapter 42's redo-only dictionary does not supply that infrastructure.

\par\medskip\noindent\begin{minipage}{\linewidth}
\begin{proposition}[Serving path and revision invariant]
Starting from an empty registry, assume only successful public transitions,
serialized calls and untampered internal dictionaries. Every active artifact
has passed offline and canary evaluation, has appeared in serving history,
and is the only artifact in active stage. Every successful switch increases
revision by one.
\end{proposition}
\begin{proof}
Creation and evaluation cannot create active stage. Promotion requires canary
with passed evidence. A switch retires the previous active before activating
its target, records that target and increments revision. Rollback requires
a retired target already in serving history with passed evaluation, so by
the induction hypothesis it originally passed canary through promotion.
No other legal transition creates active state. Revision increments once in
the common switch routine and nowhere else.
\end{proof}
\end{minipage}\par\medskip
This is a reachability invariant, not an empirical safety claim. Directly
editing the store or forging evidence violates its premises.

\section{Rollback is a governed switch, not time reversal}
\label{evod43:rollback}
\index{Rollback!eligibility}\index{state compatibility}
The incoming reference allowed rollback to any registered artifact, including
one that had never passed evaluation. We require a previously served,
validated, currently retired target and the same promoter/revision checks.
\Listing{Registry.rollback}

A \Term{state fingerprint} names the exact continuation compatibility class,
including model parameters, tokenizer, native schema and numerical semantics.
An interface version or tensor shape alone is insufficient. Incompatible
promotion or rollback requires explicit drain/reset. Existing requests should
remain pinned to their original artifact until a lifecycle fence proves their
work and output are settled. New requests start with the target's fresh or
independently validated restored state.

\Plate{03-state}{Serving identity and in-flight state have different lifetimes. The fixture's drain flag expresses a precondition that a real service must verify.}{evod43:state}
The registry stores a declared fingerprint string. It does not derive that
fingerprint from Chapter 31/35 artifacts or implement a state migration.
The coordinator must test old-request pinning, completion acknowledgement and
new-request admission against the real Chapter 30/34/35 interfaces, then the
stable Chapter 36 allocator. Local transition tests do not close that boundary.

Rollback restores future serving choice; it cannot reverse messages, trades,
deleted external objects or other irreversible actions already performed.
Memory schema rollback can also require forward-compatible readers or an
explicit migration. The active pointer, native execution state, durable
memory state and external effects are distinct coordinates.

\section{Read the lineage result and monitor after promotion}
\label{evod43:evidence}
\index{drift monitoring}
\begin{center}\input{\Generated/lineage.tex}\end{center}
The generated trace serves baseline, rejects the high-quality unsafe child,
serves compatible-by-reset v2, then returns to baseline with another reset.
The final serving revision is three. Immutable parent links and all rejected
evidence remain intact. These are deterministic fixture states, not deployment
events or measured capability improvements.

Negative controls attempt empty/missing/NaN gates, training-data evaluation,
unknown suites, wrong roles, insufficient canary observations, canary failures,
promotion before canary, stale revisions, incompatible no-drain switches and
rollback to never-served candidates. Tests verify failure leaves the serving
state unchanged where the operation is rejected before mutation.

After real promotion, monitor workload distribution, regression/safety slices,
latency/resource SLOs, memory changes and anomalous tool behavior. A drift
signal can trigger investigation, bounded rollback authority or a halt.
Unbounded automated rollback loops can themselves damage availability;
define observation windows, thresholds, cooldowns and manual escalation.
Monitoring adds evidence; it does not retrospectively justify a missing gate.

\section{Twelve worked exercises}
\label{evod43:exercises}
\begin{enumerate}
\item Why reject quality .95 and safety .5? \textbf{Answer:} Safety must
reach 1 under the fixed conjunctive policy. Quality cannot compensate.
\item Does stage change alter artifact identity? \textbf{Answer:} No.
Stage/evidence live separately from immutable content and lineage fields.
\item What prevents duplicate creation from reviving a rejected object?
\textbf{Answer:} The content ID lookup returns the existing artifact without
resetting its stage or evaluation.
\item Why are parent edges acyclic? \textbf{Answer:} A fresh child's parent
must already exist; immutable edges always point to earlier registered objects.
\item Derive the stipulated regularized optimum. \textbf{Answer:} Set
$c(\theta-b)+\lambda F(\theta-a)=0$ and divide by the positive
$c+\lambda F$, yielding the displayed weighted average.
\item Does EWC prove a compiler rewrite legal? \textbf{Answer:} No. It
regularizes parameters; data/effect observations need separate compiler tests.
\item What does zero canary errors in 20 independent trials establish?
\textbf{Answer:} Under the Bernoulli premise, a one-sided 95\% upper bound
is about 13.9\%, far from a tiny failure-rate assurance.
\item Two promoters read revision 1. What follows? \textbf{Answer:} The
first successful switch moves it to 2; the second's expected revision 1 fails.
A real concurrent implementation needs an atomic durable check/update.
\item May rollback target a rejected child? \textbf{Answer:} No. It must
be previously served, validated and currently retired.
\item Same tensor shape, different model bytes: may state carry over?
\textbf{Answer:} Not without an independently verified conversion. Exact
artifact-bound recurrence/KV semantics can change despite equal geometry.
\item Model rollback reverses external effects? \textbf{Answer:} No.
External systems need their own reconciliation or compensation, if possible.
\item Design a structural-adaptation evaluation. \textbf{Rubric:} Specify
transformation and inheritance, matched budgets, untouched final data, compiler
and native-state regressions, hard policy invariants, bounded shadow/canary
effects, authenticated promotion, state compatibility and rollback monitoring.
Synthetic registry reachability alone is not evidence of learned improvement.
\end{enumerate}

\section{Vocabulary and industrial handoff}
\Term{Lineage}: immutable derivation relations among artifacts.
\Term{Promotion}: a governed change to serving identity after evidence gates.
\Term{Rollback}: a validated switch to a previously served compatible object,
with state reset when needed. Chapter 44 turns these requirements into
packaging, deployment and observable service contracts.
